\runninghead{Note on Sources and Delivery}
\section{Note on Sources and Delivery}

\textbf{Occasion, audience, and delivery estimate.} This homily is written for an adult parish assembly for \latin{Dominica prima Adventus} (First Sunday of Advent, Sunday of the first class, \latin{Statio ad S.~Mariam maiorem}) in the 1962 typical \work{Missale Romanum} (\work{Proprium de tempore}, pp.~1--2), dated occurrence Sunday, 29~November~2026 under the universal General Roman Calendar (omitting the ordinary commemoration of St.~Saturninus and pre-1955 seasonal extra orations, and strictly isolated from postconciliar calendars, lectionaries, and prayers). The spoken body contains 1,374 spoken words (1,363 raw whitespace tokens in source), giving a word-count delivery estimate of approximately 10.6 to 11.5~minutes at an unhurried preaching pace of 120--130 words per minute (about 11.0 to 11.9~minutes at 115--125 words per minute). This figure is a word-count estimate only, not a timed human or audible delivery. The text is authored study and preaching prose, not an official liturgical book or a claim of ordained ministry.

\textbf{Relation to the reviewed studies and decisive textual detail.} The speech follows a compatible combination of the expansive and concise studies' first interpretation (\work{The World Grown Old and the Coming of the Strict Judge}) and the penitential and sacramental conversion movement of the third interpretation (\work{The Soil of Confession, the Path of Grace, and the Altar Bridge}), adding no new evidence-dependent claims. Within the first interpretation, it preserves St.~Cyril of Alexandria's promise of bodily resurrection and cosmic refashioning alongside St.~Gregory the Great's detachment, and brings to the Gospel's close St.~Augustine's teaching that the figure of this world passes away, not its nature, renewed through universal fire rather than destroyed (\work{De civitate Dei}~XX.14, 16, on Apoc.~20:11, 21:1, and 1~Cor.~7:31---not a direct commentary on Luke~21:33, which St.~Bede applies directly to Luke~21:33 in \work{In Lucae Evangelium Expositio}~VI as passing \latin{per imaginem} rather than \latin{per essentiam}, cleansed \latin{per ignem tergetur}); within the third, it retains the textual qualification that St.~Augustine's ``sweetness'' exposition of Psalm~84:13 rests on his Old Latin Psalter (\latin{suavitatem}) where the 1962 Missal sings \latin{benignitatem} (``goodness''). Linking these patristic expositions across the Gospel, Epistle, psalms, and altar prayers is the homily's editorial synthesis rather than any single Father's commentary on the complete Mass. The second interpretation's Marian-allegorical readings of the cloud (Luke~21:27) and virgin earth (Ps.~84:13) are left to the two studies so the speech sustains one continuous argument for a single hearing.

The homily opens from a familiar pastoral reaction---meeting Advent's opening prophecy of cosmic terror with evasion or anxious clinging---and turns on a decisive textual detail in Luke~21:28--30: immediately after foretelling men withering away for fear, Christ commands His disciples to lift up their heads and points to the fig tree and all the trees putting forth fruit, which St.~Gregory unlocks through the ironic reversal ``the fruit of the world is its ruin'' (balanced by St.~Cyril on bodily resurrection and St.~Augustine on fiery cosmic renewal). That Gospel turn is joined directly to the Collect's inward naming of ``approaching dangers brought on by our sins'' (contrasted with Psalm~79:3's external enemies), St.~Paul's shift from ``works of darkness'' to ``armour of light'' (Rom.~13:12, St.~John Chrysostom), and Psalm~84:8 and~13's humility and fruitful earth opened by ``the plough of confession'' (St.~Augustine), leading into the Secret, Communion, and Postcommunion.

\textbf{Exact loci, translations, and rights.} Psalm citations follow Vulgate numbering first, with Hebrew equivalents in parentheses: Introit, Ps.~24:1--3, 4b (Heb.~25:1--3a, 4); Collect, \latin{Excita, quaesumus} (echoing Ps.~79:3 [Heb.~80:3]); Epistle, Rom.~13:11--14a (omitting v.~14b); Gradual, Ps.~24:3b, 4b (Heb.~25:3a, 4); Sunday Alleluia, Ps.~84:8 (Heb.~85:8; modern English~85:7); Gospel, Luke~21:25--33; Offertory, Ps.~24:1--3 (Heb.~25:1--3a); Secret, \latin{Haec sacra nos}; Communion, Ps.~84:13 (Heb.~85:13; modern English~85:12); Postcommunion, \latin{Suscipiamus, Domine} (echoing Ps.~47:10 [Heb.~48:10; modern English~48:9]). The 1962 Latin texts were visually collated against the CMAA 1962 facsimile PDF (collation control only) and corroborated verbatim in pre-1931 public-domain Missals---Pustet (Ratisbon, 1862), Descl\'ee (Tournai, 1889), and Mame (Tours, 1922)---under 17~U.S.C.~\S103(b). Scriptural English quotes the public-domain Douay--Rheims Bible (Challoner revision, 1899 American edition, controlled against Challoner Gutenberg~1581, where Luke~21:29 reads ``a similitude'' against ``in a similitude''), and liturgical prayer wording follows the public-domain historical translation of Fr.~F.~X. Lasance and Fr.~Francis Augustine Walsh, O.S.B., \work{The New Roman Missal in Latin and English} (1945, pp.~103, 105, 106; study English only, not official 1962 liturgical English). Patristic loci used: St.~Gregory the Great, \work{Homiliae in Evangelia}, Hom.~1, \S\S3, 5--6 (PL~76, cols.~1077--1082; Wikisource Latin, CC~BY-SA~3.0; \work{Catena Aurea}, Oxford~1843); St.~Cyril of Alexandria, \work{Commentary on S.~Luke}, Sermon~CXXXIX on Luke~21:25--28 (Payne Smith, pp.~653--654); St.~Augustine of Hippo, \work{De civitate Dei}~XX.14, 16 (trans.\ M.~Dods, NPNF~I.2; applied directly to Luke~21:33 by St.~Bede, \work{In Lucae Evangelium Expositio}~VI, PL~92) and \work{Enarrationes in Psalmos} on Ps.~84, \S\S9, 14--15 (English numbering Ps.~85, \S\S6, 11--12, Coxe, NPNF~I.8, 1888); and St.~John Chrysostom, \work{Homilies on Romans}, Hom.~24 on Rom.~13:11--14 (NPNF~I.11).

\section{References}
\runninghead{References}

{\small
\begin{itemize}[itemsep=0.05em,topsep=0.2em]
\item \work{Missale Romanum ex decreto SS.~Concilii Tridentini restitutum, Summorum Pontificum cura recognitum, editio typica} (Vatican City: Typis Polyglottis Vaticanis, 1962), \work{Proprium de tempore}, \latin{Dominica prima Adventus}, pp.~1--2 (visually collated via CMAA facsimile PDF, rights unresolved, collation control only).
\item Pre-1931 public-domain Roman Missals corroborating the Latin texts under 17~U.S.C.~\S103(b): Pustet (Ratisbon, 1862), pp.~1--3; Descl\'ee (Tournai, 1889); Mame (Tours, 1922), p.~2 (Epistle).
\item Lasance, F.~X., and F.~A. Walsh, O.S.B., \work{The New Roman Missal in Latin and English} (New York: Benziger Brothers, revised 1945), pp.~103, 105, 106 (U.S.~public domain; historical study English).
\item \work{Biblia Sacra iuxta Vulgatam Clementinam} and \work{The Holy Bible Translated from the Latin Vulgate} (Douay--Rheims, Challoner revision, 1899 American edition, cross-checked against Challoner Gutenberg~1581; Ps.~24:1--4; Ps.~47:10; Ps.~79:3; Ps.~84:8, 13; Rom.~13:11--14; Luke~21:25--33; public domain).
\item Augustine of Hippo, St., \work{Enarrationes in Psalmos} on Ps.~84 (\S\S9, 14--15; Coxe Ps.~85, \S\S6, 11--12, NPNF~I.8, 1888) and \work{De civitate Dei}, Book~XX, chs.~14 and~16 (trans.\ M.~Dods, NPNF~I.2; public domain).
\item Bede, St., \work{In Lucae Evangelium Expositio}, Book~VI, on Luke~21:33 (PL~92; public domain).
\item Cyril of Alexandria, St., \work{Commentary on the Gospel according to S.~Luke}, Sermon~CXXXIX on Luke~21:25--28, trans.\ R.~Payne Smith (Oxford, 1859), pp.~653--654 (public domain).
\item Gregory the Great, St., \work{Homiliae in Evangelia}, Homily~1, \S\S3, 5--6 (PL~76, cols.~1077--1082; Wikisource Latin, CC~BY-SA~3.0; \work{Catena Aurea}, Oxford, 1843).
\item John Chrysostom, St., \work{Homilies on the Epistle to the Romans}, Homily~24 on Rom.~13:11--14 (NPNF~I.11; public domain).
\end{itemize}
}
